Sequence models of RNA-binding protein (RBP) occupancy are commonly evaluated by the area under
the receiver operating characteristic curve (AUROC) on held-out sequence windows. Positive
windows are defined by crosslinking assays such as eCLIP
\citep{vannostrand2016,vannostrand2020}; negative windows must instead be constructed. Typical
choices include unbound regions from the same transcripts \citep{maticzka2014},
composition-matched genomic windows \citep{khan2021}, and sites bound by other RBPs
\citep{horlacher2023}. This choice determines the discrimination problem presented to the
model, but it is often described only briefly.

The consequences can be substantial. \citet{horlacher2023} evaluated eleven RBP-prediction
methods across hundreds of CLIP-seq experiments and showed that apparent performance depended
on the negative-set definition. They proposed a bias-aware protocol in which sites bound by
other RBPs serve as negatives, thereby reducing some of the compositional and accessibility
differences between bound sites and genomic background. Related sensitivity has been observed
in transcription-factor binding prediction: models trained with dinucleotide-shuffled
negatives transferred less effectively to held-out high-quality data than models trained with
genomic negatives \citep{tourne2026}. More generally, evaluation-set construction can inflate
apparent accuracy, alter method rankings and change the interpretation of genomic prediction
models \citep{grimm2015,whalen2022,krutzfeldt2020,cohendavidi2024}.

Raw AUROC combines two quantities: the signal captured by a model and the difficulty of the
classification task defined by the benchmark. Uniform genomic negatives, for example, may
differ markedly from crosslinked positives in nucleotide composition. A classifier can then
achieve a high AUROC by exploiting composition alone. GC matching reduces one component of
this difference, but GC content constrains only one linear combination of the sixteen
dinucleotide frequencies. The remaining compositional degrees of freedom may still support
substantial discrimination. Similar dependence of incremental performance on baseline strength
has been described outside genomics \citep{pencina2012}.

I therefore distinguish a model's \emph{apparent AUROC} from its contribution beyond an
explicit composition baseline. Apparent AUROC is the performance reported under a particular
negative-set protocol. The nested contribution is the difference between two pooled
out-of-fold AUROCs: one from a logistic model containing composition features and the sequence
model's score, and one from the same composition features alone. Both models are fitted on the
same observations and folds. The primary baseline contains mono- and dinucleotide frequencies
and sequence entropy, represented by 19 columns after removing redundant frequency levels.
Because the resulting contribution depends on the baseline used to define it, I also evaluate
higher-order composition models.

This analysis uses 94 ENCODE eCLIP datasets and three negative-set protocols: GC-matched
genomic windows, dinucleotide-matched genomic windows, and a bias-aware set drawn from sites of
other RBPs assayed in the same cell line. The protocols differ not only in composition matching
but also in transcript-region constraints and transcriptional status. Across protocols I use
the same source peaks and hold the model class and hyperparameters, chromosome-blocked fold
assignment and estimator implementation fixed; the sequence model and composition baseline are refitted within each
protocol. I examine a 4-mer logistic regression, a 7089-parameter convolutional neural network
and a fully fine-tuned 19.78~M-parameter SpliceBERT model \citep{chen2024}.

Within the study panel, stricter composition matching lowered apparent AUROC while increasing the
estimated contribution beyond composition. The magnitude of this contribution varied several-
fold across protocols and was comparable to the variation produced by changing model class.
The dependence persisted across alternative performance measures and baseline orders, although
its magnitude changed. I further identified a positive null floor in the conventional
two-stage estimator and showed that cross-fitting the score covariate largely removed it.
Finally, an analysis of 135 held-out datasets from \citet{horlacher2023} reproduced the
protocol dependence but not the inverse relation between apparent difficulty and contribution.

Together, these results motivate three reporting practices: evaluate the composition baseline
under the same protocol as the sequence model, cross-fit score covariates when estimating
incremental performance, and avoid comparing contributions obtained under different
negative-set constructions. The aim is not to identify a universally preferred negative set,
but to make the scientific question posed by each benchmark explicit and its measurements
interpretable.
